% ===== Джерела до розділу 15. Ключі з префіксом ch15: =====
% Позиції ядра (bib-core.tex) тут НЕ дублюються.

% --- Атаки на ланцюг постачання: огляди та кейси ---
\bibitem{ch15:enisa} ENISA. Threat Landscape for Supply Chain Attacks. European Union Agency for Cybersecurity, 2021.

\bibitem{ch15:solarwinds} CISA. Advanced Persistent Threat Compromise of Government Agencies, Critical Infrastructure, and Private Sector Organizations (SolarWinds). Cybersecurity and Infrastructure Security Agency Alert AA20-352A, 2020.

\bibitem{ch15:ohm} Ohm M., Plate H., Sykosch A., Meier M. Backstabber's Knife Collection: A Review of Open Source Software Supply Chain Attacks. \emph{Detection of Intrusions and Malware, and Vulnerability Assessment (DIMVA 2020)}. LNCS Vol.~12223. Cham: Springer, 2020. P.~23--43.

\bibitem{ch15:ccleaner} Avast Threat Intelligence Team. Progress on CCleaner Investigation. Avast Blog, 2017.

% --- Цілісність та походження артефактів ---
\bibitem{ch15:slsa} Open Source Security Foundation (OpenSSF). SLSA: Supply-chain Levels for Software Artifacts. URL: \url{https://slsa.dev/}.

\bibitem{ch15:sigstore} Newman J., Torres-Arias S., et al. Sigstore: Software Signing for Everybody. \emph{Proceedings of the 2022 ACM SIGSAC Conference on Computer and Communications Security (CCS~'22)}. New York: ACM, 2022. P.~2353--2367.

\bibitem{ch15:intoto} Torres-Arias S., Afzali H., Kuppusamy T. K., Curtmola R., Cappos J. in-toto: Providing Farm-to-Table Guarantees for Bits and Bytes. \emph{Proceedings of the 28th USENIX Security Symposium}. USENIX Association, 2019. P.~1393--1410.

% --- SBOM та прозорість ---
\bibitem{ch15:sbom} National Telecommunications and Information Administration (NTIA). The Minimum Elements for a Software Bill of Materials (SBOM). U.S. Department of Commerce, 2021.

\bibitem{ch15:reproducible} Lamb C., Zacchiroli S. Reproducible Builds: Increasing the Integrity of Software Supply Chains. \emph{IEEE Software}. 2022. Vol.~39, №~2. P.~62--70.

% --- Українські джерела ---
\bibitem{ch15:supplychain-ua} Загрози ланцюга постачання програмного забезпечення та методи забезпечення цілісності й довіри до компонентів. \emph{Кібербезпека: освіта, наука, техніка}. 2024. URL: \url{https://csecurity.kubg.edu.ua/index.php/journal/article/view/802}.

\bibitem{ch15:notpetya-ua} Аналіз атаки NotPetya як прикладу компрометації ланцюга постачання програмного забезпечення в Україні. \emph{Захист інформації}. 2023. Т.~25, №~1.
